\section{The main theorem}\label{sec:main-theorem}

One family of local Wilson observables gives continuum correlations after the full ultraviolet limit at fixed volume and a refinement of any prescribed increasing volume sequence. The limiting mixture has a phase decomposition, a selected phase has a relativistic reconstruction, and one system of OPE coefficients works in every state. A degree-four field distinguishes the selected theory from every finite Wick profile of Definition~\ref{def:trace-test}.

\begin{theorem}[Existence]\label{thm:A}
For every gauge group of Definition~\ref{def:gauge-group} with simple nonabelian Lie algebra, there exists a regular Wilson family $\mathcal W$ such that, for every strictly increasing $\mathsf p:\mathbb N\to\mathbb N$, there exist a strictly increasing $\mathsf q:\mathbb N\to\mathbb N$ and hierarchy data H1–H8 over $\mathcal W$ with the following properties, for $\mathsf v=\mathsf p\circ\mathsf q$:

\begin{enumerate}[label=\textup{A-P\arabic*.},leftmargin=3.5em,nosep]
\item $\mathcal W$ converges in the sense of Definition~\ref{def:ordered-convergence} along $\mathsf v$ to
\begin{equation}\label{eq:main-kernel}
K_n(v;x)=S_n^{\mathrm{mix}}(\pi_{\rm quo}(v_1),\ldots,\pi_{\rm quo}(v_n);x)\,.
\end{equation}
\item all R1–R15 of Definition~\ref{def:reconstruction} hold.
\item O1 and O2 of Definition~\ref{def:ope} hold for the single system H7.
\item T1–T4 of Definition~\ref{def:trace-test} hold for the field H8.
\end{enumerate}
\end{theorem}

\begin{remark}
All lattice data precede the quantifier over $\mathsf p$; the refinement and every hierarchy datum follow it. Theorem~\ref{thm:A} has no additional hypotheses or conclusions implicit in the words “hierarchy”, “trace” or “existence”. In particular, it asserts neither a stress tensor nor an identification of polynomial field content. The integral convention in Definition~\ref{def:reconstruction} is part of its formulation. The relativistic theory belongs to $\xi_*$, whereas the lattice limit is $S^{\rm mix}$. The phase space need not be a point, distinct phases need not agree, and an individual phase need not arise as a limit of the lattice states. No mass gap is asserted.
\end{remark}

\begin{remark}
The volume refinement is common to every arity, label list and separated test. The ultraviolet limit uses the full sequence at every admissible fixed volume, before the volume tends to infinity. Uniqueness of the infinite-volume limit and independence of the chosen refinement are not asserted; neither an interchange of the limits nor convergence along a diagonal is claimed. If the ultraviolet index is restricted to a prescribed subsequence, the same refinement and hierarchy satisfy A-P1, since each full fixed-volume limit is also the limit along that subsequence.
\end{remark}

The construction also identifies the local field spaces, supplies a conserved stress tensor, reconstructs a theory in every phase, and gives nested products and contact distributions. The next theorem states these conclusions for every positive block-spin penalty; the principal propositions below give their quantitative content.

\begin{theorem}[The construction]\label{thm:B}
For every gauge group of Definition~\ref{def:gauge-group} with simple nonabelian Lie algebra and every real $\mathsf b>0$, there exist admissible constants, level data, local sections, reference levels and one regular Wilson family with the following properties.

\begin{enumerate}[label=\textup{B\arabic*.},leftmargin=2.7em,itemsep=0.35em,topsep=0.35em]
\item \textup{\textbf{Lattice structure.}} The block-spin transformation is Definition~\ref{def:block-spin} with penalty $\mathsf b$. The spacings and volumes satisfy Proposition~\ref{prop:U1} and Proposition~\ref{prop:U2}; the tokens are $(d,\mathbf c)$, $\mathbf c\in V_d$, and their insertions are Definition~\ref{def:insertions}. Local sections have the Taylor normalization of Definition~\ref{def:local-sections}, and all marked transfers have the reference normalization of Definition~\ref{def:insertions}. The bare couplings belong to one specified recursion of represented densities reaching $\bar\beta$ at level zero, with the exact equations and bounds of Proposition~\ref{prop:U2}. In particular,
\begin{equation}\label{eq:running-slope}
\lim_{m\to\infty}\ \sup_{N\ge m,\ p}\left|\frac{\beta_{N,p,m}}{s_m}-\kappa_G\right|=0\,.
\end{equation}

\item \textup{\textbf{Fixed volume.}} Proposition~\ref{prop:U4} holds along the full ultraviolet sequence at every volume, for every token list and every finite positional derivative order.

\end{enumerate}

For every strictly increasing volume sequence $\mathsf p$, there are a strictly increasing $\mathsf q$, a mixed table $S$ and one hierarchy over this same family, with all properties A-P1–A-P4 along $\mathsf p\circ\mathsf q$, and with the following further properties. All objects in B3–B6 refer to the table $S$, the conserved tensor $T_{ab}$ of Proposition~\ref{prop:P7}, the phase decomposition of Proposition~\ref{prop:P8} and the selected theory.

\begin{enumerate}[resume,label=\textup{B\arabic*.},leftmargin=2.7em,itemsep=0.35em,topsep=0.35em]
\item \textup{\textbf{Euclidean fields and stress.}} The fields are $F_d=V_d/\nullspace{d}(S)$ with the inclusions, derivatives and response realization in Proposition~\ref{prop:P5}, and $\dim F_d=\dim\maxwellspace{d}$. The mixed table and every phase obey E1–E7 and the $O(4)$ covariance of Proposition~\ref{prop:P7}. For its conserved tensor $T_{ab}\in F_4$ they satisfy the normalized stress laws of Definition~\ref{def:stress-laws} (Propositions~\ref{prop:P7} and~\ref{prop:P8}). The field $\Theta=\tau-S^{\xi_*}(\tau)\unitfield{4}$, $\tau=\sum_aT_{aa}$, is the selected-phase centered trace of $T$. The comparison in Proposition~\ref{prop:P7}(e), including all derivatives, arbitrary stress slots, transported derivatives and rotation generators, holds along one common scaling sequence. Its ordinary field frames are eventually exactly compatible with the structures specified in Proposition~\ref{prop:P6}.

\item \textup{\textbf{Every phase.}} The theories, continuations, surjective OS comparison maps, one-dimensional zero-energy spaces, trivial bounded weak commutants, phase-uniform low-energy and heat bounds, continued stress tensors, continuity and joint-test averages of Proposition~\ref{prop:P9} hold simultaneously for every $\xi\in\phases$.

\item \textup{\textbf{Trace and non-freeness.}} Proposition~\ref{prop:P10} holds, including the operator identification, the full-radius limit and the unique positive spectral measure with its compact-test and Schwartz-test high-energy limit. All three comparison conclusions of Proposition~\ref{prop:P11} hold for the selected theory, including its finite-profile conclusion and its localized averaged Fock countertest.

\item \textup{\textbf{Products and contacts.}} There is one OPE prescription with all separated, integrated and nested estimates of Proposition~\ref{prop:P12}, with the rates in Definition~\ref{def:nested-clusters}. With exactly the same field projections there are the real bases, subtraction scales and contact distributions of Proposition~\ref{prop:P13}.

\end{enumerate}
\end{theorem}

\begin{remark}
The gauge group and penalty are fixed first. The constants, level data, local sections, reference levels, represented trajectories and Wilson family are one choice satisfying Propositions~\ref{prop:U1} and~\ref{prop:U2}, made in the order of Section~\ref{sec:constants}. The assertion does not quantify over every choice satisfying the numerical inequalities. Only afterwards is $\mathsf p$ prescribed; one refinement serves all later conclusions. Setting $\blockpenalty=1$ and retaining A-P1--A-P4 gives Theorem~\ref{thm:A}. Theorem~\ref{thm:B} is a weakening of the construction: the degree-four position normalization of the trace and the contact estimates on general collision subspaces are proved in Part~II but are not clauses of the theorem.
\end{remark}

\begin{remark}
The selected phase $\xi_*$ is taken from the full-measure set on which the zero-energy space is one-dimensional (Proposition~\ref{prop:P8}). Proposition~\ref{prop:P9} gives this property for every phase, so the choice imposes no further condition. The contact distributions extend separated correlations; uniqueness is not asserted.
\end{remark}
